

\documentclass[notes, 11pt]{beamer} % THIS
% \documentclass[notes, 11pt]{beamer}
% \documentclass[notes, 11pt, aspectratio=169]{beamer}

% \usepackage{helvet}
% \usepackage[default]{lato}
% \usepackage{array}


% \documentclass[handout]{beamer}
% \usepackage{pgfpages}
% \pgfpagesuselayout{4 on 1}[a4paper,border shrink=5mm]

\usepackage{bbm}
\usepackage{amsthm}
\usepackage{amssymb}
\usepackage{xcolor}
\usepackage{graphicx}
\usepackage{epstopdf}
\usepackage{booktabs}
\usepackage[percent]{overpic}

\usepackage{sgamevar, tikz} % Game theory packages
\usetikzlibrary{trees, calc} % For extensive form games
\usepackage{subfig} % Manipulation and reference of small or sub figures and tables

\setbeamertemplate{theorems}[numbered] % to number
% \usetheme{Malmoe}
\usetheme{default}
% \usecolortheme{whale}
\usecolortheme{rose}


\newtheorem{defi}{Definition}
\newtheorem{teo}{Theorem}
\newtheorem{remi}{Remark}
\newtheorem{exi}{Example}
\newtheorem{propi}{Proposition}
\newtheorem{proteo}{Proof}
\newtheorem{lemi}{Lemma}
\newtheorem{prolemi}{Proof}
\newtheorem{propropi}{Proof}
\newtheorem{coro}{Corollary}
\newtheorem{proci}{Algorithm}
\newtheorem{ass}{Assumption}

\setbeamertemplate{footline}[frame number]

\usepackage{caption}
\captionsetup[figure]{labelformat=empty}%

\newenvironment{wideitemize}{\itemize\addtolength{\itemsep}{10pt}}{\enditemize}
% \usepackage{enumitem}

\usepackage[spanish]{babel}

% \renewcommand{\baselinestretch}{1}
% \expandafter\def\expandafter\insertshorttitle\expandafter{%
%   \insertshorttitle\hfill%
%   \insertframenumber\,/\,\inserttotalframenumber}

\DeclareMathOperator*{\argmax}{arg\,max}

\begin{document}
\title{Microeconom\'ia Aplicada (MAE)}   
\subtitle{Programa del Curso}   
\author[Paz y Mino ]{Jos\'e Manuel Paz y Mi\~no}
\institute[shortinst]{FEN UCHILE}
\date{Oto\~no 2025} 

\frame{\titlepage} 

\begin{frame}{Introducci\'on}
	\begin{wideitemize}
		\item<1-> Primer curso del Magister en An\'alisis Econ\'omico (MAE). \textbf{Bienvenidos}.
		\item<2-> Curso de magister distinto a pregrado.
		\begin{wideitemize}
			\item<3->[--] Mayor autonom\'ia del estudiante.
			\item<4->[--] Contenidos y evaluaciones m\'as abstractas.
			\item<5->[--] Asume el alumno tiene conocimientos previos.
		\end{wideitemize}
	\end{wideitemize}
\end{frame}

\begin{frame}{Introducci\'on}
	\begin{wideitemize}
		\item<1-> Curso MAE: Dificil de dise\~nar y dictar.
		\item<2-> Reto entre balancear rigurosidad y aplicaciones.
		\item<3-> Objetivos
		\begin{wideitemize}
			\item<4->[--] Dar al alumno confianza para que pueda consumir los principales modelos de teor\'ia de juegos existentes.
			\item<5->[--] Identificar modelos de teor\'ia de juegos en sus campos de interes \textbf{que idealmente luego guien trabajo emp\'irico y/o extensiones te\'oricas (tesis o AFE).}
		\end{wideitemize}
	\end{wideitemize}
\end{frame}

\iffalse
\begin{frame}{Introducci\'on}
	\begin{wideitemize}
		\item<1-> Teor\'ia de Juegos! Escenario competitivo donde resultado para un participante depende de acciones de otros.
		\item<2-> Ejemplos
		\begin{wideitemize}
			\item<3->[--] ¿Por qu\'e muchos paises desean tener una bomba at\'omica? (\textit{deterrence})
			\item<4->[--] ¿Por qu\'e se reglamenta la extracci\'on de recursos renovables? (\textit{bienes comunes})
			\item<5->[--] ¿Por qu\'e cachip\'un se juega de manera simult\'anea y no secuencial? (\textit{first mover advantage})
		\end{wideitemize}
	\end{wideitemize}
\end{frame}

\begin{frame}{Introducci\'on}
	\begin{wideitemize}
		\item<1-> Formalmente, Teor\'ia de Juegos es rama de matem\'atica pero con gran impacto en ciencias sociales.
		\item<2-> Nuestros jugadores son racionales.
	\end{wideitemize}
\end{frame}
\fi

\begin{frame}{Contenidos}
	\begin{wideitemize}
		\item<1-> Juegos Est\'aticos: Info completa e incompleta
		\item<2-> Juegos Din\'amicos: Info completa e incompleta
		\item<3-> Juegos Repetidos
		\item<4-> Otros: reputaci\'on, contratos, etc.
	\end{wideitemize}
\end{frame}

\begin{frame}{Evaluaciones}
	\begin{wideitemize}
		\item<1-> 1 Solemne y 1 Examen Final
		\item<2-> 2 Problem Sets
		\item<3-> 1 Proyecto de Investigaci\'on (2 entregas)
		\begin{wideitemize}
			\item<4->[--] Revisi\'on de literatura y explicaci\'on de modelo existente.
			\item<5->[--] Extensi\'on.
		\end{wideitemize}
	\end{wideitemize}
\end{frame}

\begin{frame}{Evaluaciones}
	\begin{wideitemize}
		\item<1-> ¿C\'omo estudiar para solemnes y examen final?
		\begin{wideitemize}
			\item<2->[--] Venir a clase.
			\item<3->[--] Repasar conceptos y rehacer ejercicios hechos en clase.
			\item<4->[--] Revisar ejercicios de final de cap\'itulo de textos recomendados.
			\item<5->[--] Googlear problems sets, core exams o parecidos para cursos de teoria de juegos a nivel de maestr\'ia y/o doctorado.
		\end{wideitemize}
	\end{wideitemize}
\end{frame}



\begin{frame}
\maketitle
\end{frame}

\end{document}
